% ── SECTION 5: QUANTUM PHYSICS AS CONVERGENCE, NOT APPROPRIATION (~500 words)

\section{Quantum Physics as Convergence, Not Appropriation}
\label{sec:confirmation}

% This project does not claim \Bahaullah was writing physics.
% Claim: the structural description in the writings CORRESPONDS to what
% physics discovered independently.
% Relationship: confirmation, not borrowing.
% Each tradition keeps its own vocabulary and methodology.

A clarification is essential.

I do not claim that \Bahaullah was writing physics, or that the
\bahai writings anticipate the technical formalism of quantum mechanics.
The writings are not a physics textbook.
They are spiritual revelation, addressed to questions of human meaning,
moral development, and the nature of God.

What I claim is different: that the \emph{structural description}
of reality in the \bahai writings --- existence as relational affinity, creation
as relational event --- corresponds to the relational structure that quantum
mechanics, on the relational reading, has independently arrived at through a
century of mathematics and experiment.
The third element of that description, love as the creative foundation at every
degree, has no counterpart in physics; it is where the writings, and Whitehead
in his own vocabulary, go further.

The relationship is convergence, not appropriation.
It is not confirmation either: experiments test physical theories, not the
writings, and the agreement lies in the shape of the two descriptions.
The \bahai writings do not need quantum mechanics to be true.
Quantum mechanics does not need the \bahai writings to be true.
What I document is that, independently of each other, both arrived at
the same relational structure --- which is what the two-wings principle predicts they should do.

Each tradition keeps its own vocabulary and methodology.
A \bahai scholar does not need to accept the formalism of the Relational Coherence
Framework to validate my argument.
A physicist does not need to accept the authority of the \bahai writings.
Each can assess the convergence on its own grounds.
